\section{Onsager's solution and the critical point}\label{sec:onsager}

Onsager's 1944 solution~\cite{onsager1944} of the two dimensional Ising model on the square lattice with coupling $J$ and zero field gives the free energy per site
\[
\begin{multlined}-\beta f = \log2+\frac{1}{2\pi^2}\int_0^\pi\!\!\int_0^\pi\log\Big(\cosh^2(2\beta J)\\ -\sinh(2\beta J)\big(\cos k_1+\cos k_2\big)\Big)\dd k_1\dd k_2.\end{multlined}
\]
The integrand vanishes at $k_1 = k_2 = 0$ exactly when $\cosh^2(2\beta J) = 2\sinh(2\beta J)$, that is $\sinh(2\beta_cJ) = 1$, so the critical point is $\beta_cJ = \frac12\log(1+\sqrt2)\approx0.4407$, the value predicted earlier by Kramers and Wannier~\cite{kramers1941} from the duality between high and low temperature expansions. The specific heat diverges logarithmically at $\beta_c$. The spontaneous magnetization, announced by Onsager and proved by Yang in 1952~\cite{yang1952}, is
\[
m^*(\beta) = \big(1-\sinh^{-4}(2\beta J)\big)^{1/8},\qquad\beta>\beta_c,
\]
vanishing at $\beta_c$ with exponent $\frac18$. Mean field theory predicts the exponent $\frac12$, so the two dimensional model is a rigorous counterexample to the universality of mean field exponents. The exact solution showed for the first time that a phase transition with nonclassical exponents can follow from a microscopic model by pure mathematics, a demonstration of the derivational reading of the sixth problem in equilibrium statistical mechanics.

\section{The Gibbs variational principle and variational inference}

\begin{theorem}[Gibbs variational principle]
Let $\nu$ be a probability measure on a space $\Omega$ and $H:\Omega\to\R$ measurable with $Z = \int e^{H}\dd\nu<\infty$. Then
\[
\log Z = \sup_Q\Big\{\int H\dd Q-\KL(Q\,\|\,\nu)\Big\},
\]
the supremum over probability measures $Q$ with $\int|H|\dd Q<\infty$. If in addition $\int|H|e^{H}\dd\nu<\infty$, which is equivalent to $\KL(Q^*\,\|\,\nu)<\infty$, the supremum is attained, uniquely at $Q^*(\dd\omega) = Z^{-1}e^{H(\omega)}\nu(\dd\omega)$.
\end{theorem}

\begin{proof}
Let $Q$ be a probability measure with $\int|H|\dd Q<\infty$. If $Q$ is not absolutely continuous with respect to $\nu$, then $\KL(Q\|\nu) = \infty$ and the functional being maximized is $-\infty$. Otherwise $Q\ll Q^*$ and $\frac{\dd Q}{\dd Q^*} = Ze^{-H}\frac{\dd Q}{\dd\nu}$, so that $\KL(Q\|Q^*) = \KL(Q\|\nu)-\int H\dd Q+\log Z$. Hence $\int H\dd Q-\KL(Q\|\nu) = \log Z-\KL(Q\|Q^*)\leq\log Z$, with equality if and only if $Q = Q^*$. This gives the upper bound, the uniqueness, and the attainment whenever $Q^*$ is admissible, that is, whenever $\int|H|\dd Q^* = Z^{-1}\int|H|e^H\dd\nu<\infty$. In general let $Q_n$ be $Q^*$ conditioned on $\{|H|\leq n\}$. It is admissible, and the same computation gives $\int H\dd Q_n-\KL(Q_n\|\nu) = \log Z+\log Q^*(|H|\leq n)$, which tends to $\log Z$.
\end{proof}

The theorem is the Gibbs variational principle of statistical mechanics, in the form of the variational formula for relative entropy of Donsker and Varadhan~\cite{donsker1975}. With $\nu$ the prior, $H$ the log likelihood and $Q^*$ the posterior, it says that Bayes' rule is the solution of a variational problem and that the log marginal likelihood is a free energy. Restricting $Q$ to a tractable family gives a lower bound, the evidence lower bound of variational inference, whose gap is the Kullback Leibler divergence from the family to the posterior. Restricting to product measures gives mean field theory.

For the Curie Weiss model of Section~\ref{sec:curieweiss}, take $\nu$ uniform on $\{\pm1\}^N$ and $H_N(\sigma) = \frac{\beta}{2N}\big(\sum_i\sigma_i\big)^2+h\sum_i\sigma_i$, so that $Z_N = 2^N\int e^{H_N}\dd\nu$. The product measure $Q_m$ with mean $m$ at every site has $\KL(Q_m\|\nu) = N\big(\log2-s(m)\big)$, where $s(m) = -\frac{1+m}{2}\log\frac{1+m}{2}-\frac{1-m}{2}\log\frac{1-m}{2}$ is the binary entropy, and $\int H_N\dd Q_m = N\big(\frac{\beta m^2}{2}+hm\big)+\frac\beta2(1-m^2)$. The theorem therefore gives $N^{-1}\log Z_N\geq\varphi(m)$ for every $m\in[-1,1]$, where $\varphi(m) = \frac{\beta m^2}{2}+hm+s(m)$. This is not the function $-\Psi(m)$ of the Hubbard Stratonovich computation. Since $\log(2\cosh u) = \max_{m'}\{um'+s(m')\}$, attained at $m' = \tanh u$, we have $-\Psi(m) = \max_{m'}\{(\beta m+h)m'+s(m')\}-\frac{\beta m^2}{2}$. Taking $m' = m$ gives $-\Psi(m)\geq\varphi(m)$, with equality if and only if $m = \tanh(\beta m+h)$. Maximizing over $m$ first, which gives $\frac{\beta m'^2}{2}$, shows that $\max_m(-\Psi) = \max_{m'}\varphi$. The two functions thus have the same stationary points, the solutions of the mean field equation, and the same maximum, which is the limit of $N^{-1}\log Z_N$. For the Curie Weiss model the product bound is therefore sharp. Mean field theory in physics and variational Bayes in statistics are the same approximation, with the same failure mode: a product family cannot represent the correlations that dominate near a phase transition.

\section{Belief propagation and the Bethe approximation}

On a tree, marginals of a Gibbs measure can be computed exactly by passing messages along edges: the message from $x$ to $y$ summarizes the effect of the subtree behind $x$, and marginals are products of incoming messages. For the Ising model with fields $h_x$ and couplings $J_{xy}$, writing messages as effective fields $u_{x\to y}$, the update is
\[
u_{x\to y} = \operatorname{artanh}\Big(\tanh(\beta J_{xy})\tanh\Big(h_x+\sum_{z\sim x,z\neq y}u_{z\to x}\Big)\Big).
\]
On graphs with cycles the same iteration, loopy belief propagation, computes stationary points of the Bethe free energy~\cite{yedidia2005}, a variational approximation that is exact on trees and accurate on locally tree like graphs. Pearl introduced belief propagation for Bayesian networks in 1982~\cite{pearl1982}, Gallager had used it for decoding in 1963~\cite{gallager1963}, and the cavity method of M\'ezard, Parisi and Virasoro is its replica symmetric physics~\cite{mezard2009}. Approximate message passing~\cite{donoho2009,bayati2011}, used earlier in this chapter for the planted clique, is its dense graph limit, and its state evolution, a one dimensional recursion for the effective signal to noise ratio, is the mechanism by which algorithmic thresholds such as $\sqrt{n/e}$ are computed.

\section*{Exercises}
\begin{enumerate}
\item Solve the one dimensional Ising model on a ring by the transfer matrix method and show that there is no phase transition at positive temperature.
\item Verify directly that the zeros of the partition function of a single Ising bond in a uniform field lie on the unit circle.
\item Show that the Curie Weiss mean field equation with $h = 0$ and $\beta>1$ has three solutions, and determine which minimize $\Psi$.
\item Prove that the Edwards Sokal measure has the Ising measure as its spin marginal.
\item Show that ridge regression with penalty $\alpha$, the minimizer of $n^{-1}\|y-X\theta\|^2+\alpha\|\theta\|^2$, is the posterior mean under the prior $N(0,(\sigma^2/(n\alpha))I)$ and Gaussian noise of variance $\sigma^2$.
\end{enumerate}
